\chapter{Foundations and main statements}
\label{chap:foundations}

% archon:covers MR4736527GeometricLocalSystemsOnVeryGeneralCurvesAndIsomonodromy/Basic.lean

This chapter collects the foundational definitions of \S1 of Landesman--Litt and
states the main Hodge-theoretic results of the paper, culminating in the project
goal, \cref{thm:very-general-VHS}. The objects ``flat vector bundle with regular
singularities'', ``isomonodromic deformation'', and ``nearby curve'' are
developed in the chapter on isomonodromic deformations, and ``polarizable complex
variation of Hodge structure'' in the Hodge chapter; here we only refer to them.

\section{Hyperbolic pointed curves}

\begin{definition}[Hyperbolic pointed curve]
  \label{def:hyperbolic}
  \lean{LandesmanLitt.Hyperbolic}
  % SOURCE: [Landesman--Litt], Definition 1.1 (definition:hyperbolic) (read from references/MR4736527-local-systems-isomonodromy.tex)
  % SOURCE QUOTE: "Let $C$ be a curve over $\mathbb C$ of genus $g$ and $D \subset C$ a
  % reduced effective divisor of degree $n$.
  % Call $(C, D)$ {\em hyperbolic} if $C$ is a smooth proper connected curve
  % and
  % either $g \geq 2$ and $n\geq 0$, $g = 1$ and $n > 0$, or $g =0$ and $n > 2$.
  % We call an $n$-pointed curve $(C, x_1, \ldots, x_n)$ {\em hyperbolic} if
  % $(C, x_1 +\cdots + x_n)$ is hyperbolic."
  \textit{Source: Landesman--Litt, Definition 1.1.}
  Let $C$ be a smooth proper connected curve over $\C$ of genus $g$, and let
  $D \subset C$ be a reduced effective divisor of degree $n$. The pair $(C, D)$
  is \emph{hyperbolic} if one of the following holds:
  \[
    g \ge 2 \text{ and } n \ge 0, \qquad
    g = 1 \text{ and } n > 0, \qquad
    g = 0 \text{ and } n > 2.
  \]
  An $n$-pointed curve $(C, x_1, \ldots, x_n)$ is \emph{hyperbolic} if
  $(C, x_1 + \cdots + x_n)$ is hyperbolic.

  % SOURCE: [Landesman--Litt], Remark following Definition 1.1 (read from references/MR4736527-local-systems-isomonodromy.tex)
  % SOURCE QUOTE: "Equivalently, $(C,D)$ is hyperbolic if and only if it has no infinitesimal automorphisms,
  % i.e., $H^0(C, T_C(-D)) = 0$."
  Equivalently, $(C, D)$ is hyperbolic if and only if it has no infinitesimal
  automorphisms, i.e.\ $H^0(C, T_C(-D)) = 0$.
\end{definition}

\section{Analytically general and very general points}

\begin{definition}[Analytically general point]
  \label{def:analytically-general}
  \lean{LandesmanLitt.AnalyticallyGeneral}
  \uses{def:moduli-stack}
  % SOURCE: [Landesman--Litt], Definition 1.2 (definition:general) (read from references/MR4736527-local-systems-isomonodromy.tex)
  % SOURCE QUOTE: "A property holds for an {\em analytically general} point of a complex orbifold $X$,
  % if there exists a nowhere dense
  % closed analytic subset $S \subset X$ so that the property holds
  % on $X - S$.
  % We say that a property holds for an
  % {\em analytically very general} point if, locally on $X$, there exists a countable collection of nowhere dense closed analytic subsets such that the property holds on the complement of their union. If $\mathscr{M}_{g,n}$ is the analytic moduli stack of
  % $n$-pointed curves of genus $g$, we say that a property holds for an
  % analytically (very) general $n$-pointed curve if it holds for an analytically
  % (very) general point of $\mathscr{M}_{g,n}$."
  \textit{Source: Landesman--Litt, Definition 1.2.}
  A property holds for an \emph{analytically general} point of a complex orbifold
  $X$ if there exists a nowhere dense closed analytic subset $S \subset X$ such
  that the property holds on $X - S$. Equivalently (and more in line with the very
  general notion below), the property holds analytically generally if, locally on
  $X$, there is a nowhere dense closed analytic subset on whose complement the
  property holds.
\end{definition}

\begin{definition}[Analytically very general point]
  \label{def:analytically-very-general}
  \lean{LandesmanLitt.AnalyticallyVeryGeneral}
  \uses{def:moduli-stack, def:analytically-general}
  % SOURCE: [Landesman--Litt], Definition 1.2 (definition:general) (read from references/MR4736527-local-systems-isomonodromy.tex)
  % SOURCE QUOTE: "We say that a property holds for an
  % {\em analytically very general} point if, locally on $X$, there exists a countable collection of nowhere dense closed analytic subsets such that the property holds on the complement of their union. If $\mathscr{M}_{g,n}$ is the analytic moduli stack of
  % $n$-pointed curves of genus $g$, we say that a property holds for an
  % analytically (very) general $n$-pointed curve if it holds for an analytically
  % (very) general point of $\mathscr{M}_{g,n}$."
  \textit{Source: Landesman--Litt, Definition 1.2.}
  A property holds for an \emph{analytically very general} point of a complex
  orbifold $X$ if, locally on $X$, there exists a countable collection of nowhere
  dense closed analytic subsets such that the property holds on the complement of
  their union. A property holds for an analytically (very) general $n$-pointed
  curve of genus $g$ if it holds for an analytically (very) general point of the
  analytic moduli stack $\sM_{g,n}$ (\cref{def:moduli-stack}).
\end{definition}

\begin{definition}[Analytic moduli stack of pointed curves]
  \label{def:moduli-stack}
  \lean{LandesmanLitt.ModuliStack}
  % SOURCE: [Landesman--Litt], \S1 (read from references/MR4736527-local-systems-isomonodromy.tex)
  % SOURCE QUOTE: "We use
  % $\mathscr{M}_{g,n}$ to denote the analytic moduli stack of smooth proper curves with
  % geometrically connected fibers and $n$ distinct marked points."
  \textit{Source: Landesman--Litt, \S1.}
  We write $\sM_{g,n}$ for the analytic moduli stack of smooth proper curves of
  genus $g$ with geometrically connected fibers and $n$ distinct marked points.
  For the purposes of this formalization, $\sM_{g,n}$ is modeled as an opaque
  complex orbifold, equipped with the analytically general and analytically very
  general predicates of \cref{def:analytically-general,def:analytically-very-general}.
\end{definition}

\section{Local systems and monodromy}

\begin{definition}[Complex local system]
  \label{def:local-system}
  \lean{LandesmanLitt.LocalSystem}
  % SOURCE: [Landesman--Litt], \S1 (corollary:non-density-of-geometric-local-systems context) (read from references/MR4736527-local-systems-isomonodromy.tex)
  % SOURCE QUOTE: "$\on{Hom}(\pi_1(C\setminus\{x_1, \cdots, x_n\}), \on{GL}_r(\mathbb{C}))\sslash \on{GL}_r(\mathbb{C})$ be the \emph{character variety} parametrizing conjugacy classes of semisimple representations of $\pi_1(C\setminus \{x_1, \cdots, x_n\})$ into $\on{GL}_r(\mathbb{C})$."
  \textit{Source: Landesman--Litt, \S1.}
  A rank-$r$ \emph{complex local system} $\bV$ on a connected complex manifold or
  punctured curve $X$ is a locally constant sheaf of $r$-dimensional $\C$-vector
  spaces on $X$. Fixing a base point, this is the same datum as a representation
  $\pi_1(X) \to \GL_r(\C)$, well defined up to conjugacy; we write
  $\rk \bV = r$ for its rank.
\end{definition}

\begin{definition}[Monodromy representation]
  \label{def:monodromy}
  \lean{LandesmanLitt.Monodromy}
  \uses{def:local-system}
  % SOURCE: [Landesman--Litt], \S1 (definition of unitary monodromy) (read from references/MR4736527-local-systems-isomonodromy.tex)
  % SOURCE QUOTE: "the associated monodromy representation
  % $\rho:\pi_1(C)\to \on{GL}_n(\mathbb{C})$"
  \textit{Source: Landesman--Litt, \S1.}
  The \emph{monodromy representation} of a rank-$n$ local system $\bV$ on $X$ (or
  of a flat vector bundle, via its underlying local system of flat sections) is
  the representation
  \[
    \rho : \pi_1(X) \longrightarrow \GL_n(\C)
  \]
  obtained by parallel transport. Its image $\rho(\pi_1(X)) \subseteq \GL_n(\C)$
  is the \emph{monodromy group}.
\end{definition}

\begin{definition}[Infinite monodromy]
  \label{def:infinite-monodromy}
  \lean{LandesmanLitt.InfiniteMonodromy}
  \uses{def:monodromy}
  % SOURCE: [Landesman--Litt], Theorem 1.3 (theorem:very-general-VHS) (read from references/MR4736527-local-systems-isomonodromy.tex)
  % SOURCE QUOTE: "$\mathbb{V}$ is a $\mathscr{O}_K$-local system on $C\setminus\{x_1, \cdots, x_n\}$ with infinite monodromy."
  \textit{Source: Landesman--Litt, \S1.}
  A local system $\bV$ (or flat vector bundle) has \emph{infinite monodromy} if
  its monodromy group $\rho(\pi_1(X))$ is infinite.
\end{definition}

\begin{definition}[Unitary monodromy]
  \label{def:unitary-monodromy}
  \lean{LandesmanLitt.UnitaryMonodromy}
  \uses{def:monodromy}
  % SOURCE: [Landesman--Litt], \S1 (read from references/MR4736527-local-systems-isomonodromy.tex)
  % SOURCE QUOTE: "In what follows, we say a flat vector bundle has \emph{unitary monodromy} if the associated monodromy representation
  % $\rho:\pi_1(C)\to \on{GL}_n(\mathbb{C})$
  % has image with compact closure."
  \textit{Source: Landesman--Litt, \S1.}
  A flat vector bundle (equivalently, its local system of flat sections) has
  \emph{unitary monodromy} if the associated monodromy representation
  $\rho : \pi_1(C) \to \GL_n(\C)$ has image with compact closure.
\end{definition}

\begin{definition}[$\sO_K$-local system]
  \label{def:OK-local-system}
  \lean{LandesmanLitt.OKLocalSystem}
  \uses{def:local-system}
  % SOURCE: [Landesman--Litt], Theorem 1.3 (theorem:very-general-VHS) (read from references/MR4736527-local-systems-isomonodromy.tex)
  % SOURCE QUOTE: "Let $K$ be a number field with ring of integers
  % $\mathscr{O}_K$. Suppose
  % $(C, x_1, \cdots, x_n)$
  % is an analytically very general $n$-pointed
  % hyperbolic curve of genus $g$, and $\mathbb{V}$ is a $\mathscr{O}_K$-local system on $C\setminus\{x_1, \cdots, x_n\}$ with infinite monodromy."
  \textit{Source: Landesman--Litt, Theorem 1.3.}
  Let $K$ be a number field with ring of integers $\sO_K$. An \emph{$\sO_K$-local
  system} on a space $X$ is a locally constant sheaf of $\sO_K$-modules on $X$
  that is locally free of finite rank; its rank is written $\rk_{\sO_K}(\bV)$.
  For each embedding $\iota : \sO_K \to \C$ it gives rise to a complex local
  system $\bV \otimes_{\sO_K, \iota} \C$ of rank $\rk_{\sO_K}(\bV)$.
\end{definition}

\begin{definition}[Local system of geometric origin]
  \label{def:geometric-origin}
  \lean{LandesmanLitt.GeometricOrigin}
  \uses{def:local-system}
  % SOURCE: [Landesman--Litt], \S1 (corollary:geometric-local-systems context) (read from references/MR4736527-local-systems-isomonodromy.tex)
  % SOURCE QUOTE: "Let $X$ be a smooth variety. We say a complex local system $\mathbb{V}$ on $X$ is \emph{of geometric origin} if there exists a dense open $U\subset X$, and a smooth proper morphism $f: Y\to U$ such that $\mathbb{V}|_U$ is a direct summand of $R^if_*\mathbb{C}$ for some $i\geq 0$."
  \textit{Source: Landesman--Litt, \S1.}
  Let $X$ be a smooth variety. A complex local system $\bV$ on $X$ is \emph{of
  geometric origin} if there exists a dense open subset $U \subseteq X$ and a
  smooth proper morphism $f : Y \to U$ such that $\bV|_U$ is a direct summand of
  $R^i f_* \C$ for some $i \ge 0$.
\end{definition}

\begin{definition}[Character variety]
  \label{def:character-variety}
  \lean{LandesmanLitt.CharacterVariety}
  \uses{def:local-system, def:monodromy}
  % SOURCE: [Landesman--Litt], \S1 (corollary:non-density-of-geometric-local-systems context) (read from references/MR4736527-local-systems-isomonodromy.tex)
  % SOURCE QUOTE: "$$\mathscr{M}_{B,r}(C\setminus \{x_1, \cdots, x_n\}):=\on{Hom}(\pi_1(C\setminus\{x_1, \cdots, x_n\}), \on{GL}_r(\mathbb{C}))\sslash \on{GL}_r(\mathbb{C})$$ be the \emph{character variety} parametrizing conjugacy classes of semisimple representations of $\pi_1(C\setminus \{x_1, \cdots, x_n\})$ into $\on{GL}_r(\mathbb{C})$."
  \textit{Source: Landesman--Litt, \S1.}
  For a punctured curve $C \setminus \{x_1, \ldots, x_n\}$, the \emph{character
  variety} is the geometric invariant theory quotient
  \[
    \sM_{B,r}(C \setminus \{x_1, \ldots, x_n\})
      := \Hom\!\bigl(\pi_1(C \setminus \{x_1, \ldots, x_n\}), \GL_r(\C)\bigr)
         /\!/ \GL_r(\C),
  \]
  parametrizing conjugacy classes of semisimple representations of
  $\pi_1(C \setminus \{x_1, \ldots, x_n\})$ into $\GL_r(\C)$.
\end{definition}

\section{Main Hodge-theoretic statements}

\begin{theorem}[Rank bound for very general VHS]
  \label{thm:very-general-VHS}
  \lean{LandesmanLitt.very_general_VHS}
  \uses{thm:isomonodromic-deformation-CVHS, lem:integral-and-unitary-implies-finite,
        def:hyperbolic, def:analytically-very-general, def:OK-local-system,
        def:infinite-monodromy, def:complex-variation, def:flat-bundle}
  % SOURCE: [Landesman--Litt], Theorem 1.3 (theorem:very-general-VHS) (read from references/MR4736527-local-systems-isomonodromy.tex)
  % SOURCE QUOTE: "Let $K$ be a number field with ring of integers
  % $\mathscr{O}_K$. Suppose
  % $(C, x_1, \cdots, x_n)$
  % is an analytically very general $n$-pointed
  % hyperbolic curve of genus $g$, and $\mathbb{V}$ is a $\mathscr{O}_K$-local system on $C\setminus\{x_1, \cdots, x_n\}$ with infinite monodromy.
  % Suppose additionally that for each embedding $\iota: \mathscr{O}_K\to \mathbb{C}$, $\mathbb{V}\otimes_{\mathscr{O}_K, \iota}\mathbb{C}$ underlies a polarizable complex variation of Hodge structure.
  % Then, $$\on{rk}_{\mathscr{O}_K}(\mathbb{V})\geq 2\sqrt{g+1}.$$"
  \textit{Source: Landesman--Litt, Theorem 1.3.}
  Let $K$ be a number field with ring of integers $\sO_K$. Suppose
  $(C, x_1, \ldots, x_n)$ is an analytically very general $n$-pointed hyperbolic
  curve of genus $g$, and $\bV$ is an $\sO_K$-local system on
  $C \setminus \{x_1, \ldots, x_n\}$ with infinite monodromy. Suppose
  additionally that for each embedding $\iota : \sO_K \to \C$, the complex local
  system $\bV \otimes_{\sO_K, \iota} \C$ underlies a polarizable complex
  variation of Hodge structure. Then
  \[
    \rk_{\sO_K}(\bV) \ge 2\sqrt{g+1}.
  \]
\end{theorem}

% SOURCE QUOTE PROOF: "We will deduce the above results from
% \autoref{theorem:isomonodromic-deformation-CVHS} below,
% using that a discrete subset of the image of a unitary $\rho$ is finite."
\begin{proof}
  \uses{thm:isomonodromic-deformation-CVHS, lem:integral-and-unitary-implies-finite,
        def:unitary-monodromy, def:infinite-monodromy}
  We argue by contradiction: suppose $\rk_{\sO_K}(\bV) < 2\sqrt{g+1}$. Fix an
  embedding $\iota : \sO_K \to \C$ and let $(E, \nabla)$ be the flat vector bundle
  with regular singularities underlying $\bV \otimes_{\sO_K, \iota} \C$, so that
  $\rk E = \rk_{\sO_K}(\bV) < 2\sqrt{g+1}$. Because $(C, x_1, \ldots, x_n)$ is
  analytically very general, an isomonodromic deformation of $(E, \nabla)$ to an
  analytically general nearby pointed curve still underlies a polarizable complex
  variation of Hodge structure. By \cref{thm:isomonodromic-deformation-CVHS},
  $(E, \nabla)$ then has unitary monodromy.

  Thus the monodromy representation $\rho$ has image with compact closure. On the
  other hand, the $\sO_K$-structure forces the image of $\rho$ to be a discrete
  subset of $\GL_n(\C)$: a discrete subset of a set with compact closure is
  finite, so the monodromy group is finite. This is recorded in
  \cref{lem:integral-and-unitary-implies-finite}. A finite monodromy group
  contradicts the hypothesis that $\bV$ has infinite monodromy. Hence
  $\rk_{\sO_K}(\bV) \ge 2\sqrt{g+1}$.
\end{proof}

\begin{theorem}[Isomonodromic deformations underlying CVHS are unitary]
  \label{thm:isomonodromic-deformation-CVHS}
  \lean{LandesmanLitt.isomonodromic_deformation_CVHS}
  \uses{def:hyperbolic, def:isomonodromy, def:nearby, def:complex-variation,
        def:unitary-monodromy, def:flat-bundle}
  % SOURCE: [Landesman--Litt], Theorem 1.6 (theorem:isomonodromic-deformation-CVHS) (read from references/MR4736527-local-systems-isomonodromy.tex)
  % SOURCE QUOTE: "Let $(C, x_1, \cdots, x_n)$ be an $n$-pointed hyperbolic curve of genus $g$.
  % Let $({E}, \nabla)$ be a flat vector bundle on
  % $C$ with $\on{rk}{E}<2\sqrt{g+1}$ and with regular singularities at the $x_i$. If an isomonodromic
  % deformation of ${(E, \nabla)}$ to an analytically general nearby $n$-pointed curve underlies a polarizable complex variation of Hodge structure, then $({E},\nabla)$ has unitary monodromy."
  \textit{Source: Landesman--Litt, Theorem 1.6.}
  Let $(C, x_1, \ldots, x_n)$ be an $n$-pointed hyperbolic curve of genus $g$.
  Let $(E, \nabla)$ be a flat vector bundle on $C$ with $\rk E < 2\sqrt{g+1}$ and
  with regular singularities at the $x_i$. If an isomonodromic deformation of
  $(E, \nabla)$ to an analytically general nearby $n$-pointed curve underlies a
  polarizable complex variation of Hodge structure, then $(E, \nabla)$ has
  unitary monodromy.
\end{theorem}

\begin{proof}
  \uses{lem:unitary, cor:unstable, cor:stable-parabolic}
  The proof is carried out in the chapter on Hodge-theoretic stability results.
  In outline: by the Harder--Narasimhan analysis, the small-rank hypothesis
  $\rk E < 2\sqrt{g+1}$ forces an isomonodromic deformation of $(E, \nabla)$ to an
  analytically general nearby curve to be parabolically semistable
  (\cref{cor:stable-parabolic}). A flat vector bundle underlying a polarizable
  complex variation of Hodge structure whose underlying bundle is semistable must
  have unitary monodromy, because a non-unitary variation produces an unstable
  Hodge filtration (\cref{cor:unstable}); this contrast is packaged as
  \cref{lem:unitary}. Combining these, $(E, \nabla)$ has unitary monodromy.
\end{proof}

\begin{corollary}[Geometric local systems on very general curves]
  \label{cor:geometric-local-systems}
  \lean{LandesmanLitt.geometric_local_systems}
  \uses{thm:very-general-VHS, def:geometric-origin, def:infinite-monodromy}
  % SOURCE: [Landesman--Litt], Corollary 1.4 (corollary:geometric-local-systems) (read from references/MR4736527-local-systems-isomonodromy.tex)
  % SOURCE QUOTE: "Let $(C, x_1, \cdots, x_n)$ be an analytically very general hyperbolic $n$-pointed curve of genus $g$. If $\mathbb{V}$ is a local system on $C\setminus\{x_1, \cdots, x_n\}$ of geometric origin and with infinite monodromy, then $\dim_{\mathbb{C}}\mathbb{V}\geq 2\sqrt{g+1}$."
  \textit{Source: Landesman--Litt, Corollary 1.4.}
  Let $(C, x_1, \ldots, x_n)$ be an analytically very general hyperbolic
  $n$-pointed curve of genus $g$. If $\bV$ is a local system on
  $C \setminus \{x_1, \ldots, x_n\}$ of geometric origin and with infinite
  monodromy, then $\dim_{\C} \bV \ge 2\sqrt{g+1}$.
\end{corollary}

% SOURCE QUOTE PROOF: "As local systems of geometric origin satisfy the hypotheses of \autoref{theorem:very-general-VHS}, we have:"
\begin{proof}
  \uses{thm:very-general-VHS, def:geometric-origin}
  A local system of geometric origin is a direct summand of $R^i f_* \C$ for a
  smooth proper morphism $f$; by the theory of polarizable variations of Hodge
  structure such a summand carries a polarizable complex variation of Hodge
  structure, and it is integral (indeed it underlies a $\Z$-, hence $\sO_K$-, local
  system). Thus $\bV$ satisfies the hypotheses of \cref{thm:very-general-VHS} with
  $\sO_K = \Z$, and that theorem gives
  $\dim_{\C} \bV = \rk_{\Z}(\bV) \ge 2\sqrt{g+1}$.
\end{proof}

\begin{corollary}[Abelian schemes over very general curves]
  \label{cor:abelian-schemes}
  \lean{LandesmanLitt.abelian_schemes}
  \uses{cor:geometric-local-systems}
  % SOURCE: [Landesman--Litt], Corollary 1.5 (corollary:abelian-schemes) (read from references/MR4736527-local-systems-isomonodromy.tex)
  % SOURCE QUOTE: "If $(C, x_1, \ldots, x_n)$ is an analytically
  % very general hyperbolic $n$-pointed genus $g$ curve, then any non-isotrivial
  % abelian scheme over $C\setminus\{x_1, \cdots, x_n\}$  has relative dimension at least $\sqrt{g+1}$.
  % Similarly, any relative smooth proper curve over $C\setminus\{x_1, \cdots, x_n\}$
  % has genus at least $\sqrt{g+1}$."
  \textit{Source: Landesman--Litt, Corollary 1.5.}
  If $(C, x_1, \ldots, x_n)$ is an analytically very general hyperbolic
  $n$-pointed genus $g$ curve, then any non-isotrivial abelian scheme over
  $C \setminus \{x_1, \ldots, x_n\}$ has relative dimension at least
  $\sqrt{g+1}$. Similarly, any relative smooth proper curve over
  $C \setminus \{x_1, \ldots, x_n\}$ has genus at least $\sqrt{g+1}$.
\end{corollary}

\begin{proof}
  \uses{cor:geometric-local-systems}
  A non-isotrivial abelian scheme $A \to C \setminus \{x_1, \ldots, x_n\}$ of
  relative dimension $d$ gives a geometric local system $R^1 \pi_* \C$ of rank
  $2d$ with infinite monodromy (non-isotriviality forces the monodromy to be
  infinite). By \cref{cor:geometric-local-systems}, $2d \ge 2\sqrt{g+1}$, i.e.\
  $d \ge \sqrt{g+1}$. For a relative smooth proper curve of genus $h$, its
  Jacobian is an abelian scheme of relative dimension $h$, and the same argument
  gives $h \ge \sqrt{g+1}$.
\end{proof}

\begin{proposition}[Maps to Hilbert modular stacks]
  \label{proposition:hilbert-modular}
  \lean{LandesmanLitt.hilbert_modular}
  \uses{thm:very-general-VHS, cor:abelian-schemes, def:hyperbolic,
        def:analytically-very-general, def:infinite-monodromy}
  % SOURCE: [Landesman--Litt], Proposition (proposition:hilbert-modular), \S7.4 (read from references/MR4736527-local-systems-isomonodromy.tex)
  % SOURCE QUOTE: "Let $\mathscr X_K$ be the Hilbert modular stack associated to a totally
  % real field $K$.
  % Let $(C, x_1, \ldots, x_n)$ be an analytically very general hyperbolic $n$-pointed
  % curve of genus $g \geq 1$.
  % Any map $\phi: C\setminus \{x_1, \ldots, x_n\} \to \mathscr X_K$ is constant."
  \textit{Source: Landesman--Litt, \S7.4.}
  Let $\mathscr{X}_K$ be the Hilbert modular stack associated to a totally real field
  $K$ (parametrizing $h$-dimensional principally polarized abelian varieties $A$
  with an injection $\sO_K \hookrightarrow \on{End}(A)$, where $h = [K:\Q]$). Let
  $(C, x_1, \ldots, x_n)$ be an analytically very general hyperbolic $n$-pointed
  curve of genus $g \ge 1$. Then any map
  $\phi : C \setminus \{x_1, \ldots, x_n\} \to \mathscr{X}_K$ is constant.
\end{proposition}

% SOURCE QUOTE PROOF: "We first use the map $\phi$ to produce a rank $2$ integral local system
% on $C \setminus \{x_1, \ldots, x_n\}$. ... after passing to a finite extension $K'$ of $K$,
% ... $R^1 f_* \mathbb Z \otimes_{\mathscr O_K} \mathscr O_{K'}$ is an $\mathscr O_{K'}$-local system of rank $2$.
% ... Hence by \autoref{theorem:very-general-VHS}, $R^1 f_* \mathbb Z\otimes_{\mathscr O_K} \mathscr O_{K'}$ has finite monodromy, as
% $2 = \rk_{\mathscr O_K'}R^1 f_* \mathbb Z\otimes_{\mathscr O_K} \mathscr O_{K'} < 2 \sqrt{g+1}$.
% Hence, the same holds for $R^1 f_* \mathbb Z$. This implies the map $\phi$ is constant using the theorem
% of the fixed part \cite[Corollaire 4.1.2]{deligne:hodge-ii} (see the proof of \autoref{corollary:abelian-schemes} for a similar argument)."
\begin{proof}
  \uses{thm:very-general-VHS, cor:abelian-schemes}
  A map $\phi : C \setminus \{x_1, \ldots, x_n\} \to \mathscr{X}_K$ yields an abelian
  scheme $f : A \to C \setminus \{x_1, \ldots, x_n\}$ with an $\sO_K$-action, so
  that $R^1 f_* \Z$ is a complex PVHS of rank $2h$ carrying an $\sO_K$-action.
  After passing to a finite extension $K'$ of $K$ (for instance the Hilbert class
  field), every rank-$2$ locally free $\sO_K$-module becomes free, so
  $R^1 f_* \Z \otimes_{\sO_K} \sO_{K'}$ is an $\sO_{K'}$-local system of rank $2$.
  For each embedding $\sO_{K'} \to \C$, this local system underlies a $\C$-PVHS.
  Since $g \ge 1$ gives $2 = \rk_{\sO_{K'}}\bigl(R^1 f_* \Z \otimes_{\sO_K}
  \sO_{K'}\bigr) < 2\sqrt{g+1}$, \cref{thm:very-general-VHS} shows
  $R^1 f_* \Z \otimes_{\sO_K} \sO_{K'}$ has finite monodromy, hence so does
  $R^1 f_* \Z$. By the theorem of the fixed part
  (\cite[Corollaire 4.1.2]{deligne:hodge-ii}; compare the proof of
  \cref{cor:abelian-schemes}), $\phi$ is constant.
\end{proof}

\begin{corollary}[Non-density of geometric local systems]
  \label{cor:non-density}
  \lean{LandesmanLitt.non_density}
  \uses{cor:geometric-local-systems, def:character-variety, lem:finite-image-not-dense}
  % SOURCE: [Landesman--Litt], Corollary 1.7 (corollary:non-density-of-geometric-local-systems) (read from references/MR4736527-local-systems-isomonodromy.tex)
  % SOURCE QUOTE: "Let $(C, x_1, \cdots, x_n)$ be an analytically very general hyperbolic $n$-pointed curve of genus $g$. Then if $1<r<2\sqrt{g+1}$, the local systems of geometric origin are not Zariski-dense in the character variety $\mathscr{M}_{B,r}(C\setminus \{x_1, \cdots, x_n\})$."
  \textit{Source: Landesman--Litt, Corollary 1.7.}
  Let $(C, x_1, \ldots, x_n)$ be an analytically very general hyperbolic
  $n$-pointed curve of genus $g$. If $1 < r < 2\sqrt{g+1}$, then the local
  systems of geometric origin are not Zariski-dense in the character variety
  $\sM_{B,r}(C \setminus \{x_1, \ldots, x_n\})$.
\end{corollary}

\begin{proof}
  \uses{cor:geometric-local-systems, def:character-variety, lem:finite-image-not-dense}
  By \cref{cor:geometric-local-systems}, every rank-$r$ local system of geometric
  origin with $r < 2\sqrt{g+1}$ has finite monodromy. Points of the character
  variety corresponding to representations with finite image are not Zariski-dense
  when $r > 1$ (\cref{lem:finite-image-not-dense}): they lie in a proper closed
  subset. Hence the geometric local systems, all of which have finite monodromy in
  this rank range, fail to be Zariski-dense in
  $\sM_{B,r}(C \setminus \{x_1, \ldots, x_n\})$.
\end{proof}
